% =====================================================================
%  CAPITULO 7 - ANALISE NODAL
% =====================================================================
\section{Análise Nodal}

\begin{conceito}[Antes de começar]
A análise nodal usa a \textbf{lei dos nós} de Kirchhoff para achar os potenciais
dos nós. Roteiro:
\begin{enumerate}[leftmargin=1.7em,itemsep=1pt,topsep=2pt]
\item Escolha um nó como \textbf{referência} (terra, $V=0$).
\item Nomeie os potenciais dos demais nós ($V_1, V_2, \dots$).
\item Em cada nó, escreva que a soma das correntes que \emph{saem} é nula. A
      corrente por um resistor entre os nós $a$ e $b$ é
      $\dfrac{V_a - V_b}{R}$.
\item Resolva o sistema.
\end{enumerate}
Uma fonte de corrente que injeta $I$ no nó entra como termo independente; uma
fonte de tensão entre dois nós cria uma relação direta entre eles.
\end{conceito}

\legendaniveis

\blocofixacao

\begin{questao}[1]{}
Na análise nodal, o nó de referência (terra) é escolhido para:
\begin{alternativas}
\item ter o maior potencial do circuito.
\item servir como ponto de potencial zero, ao qual os demais são comparados.
\item ser sempre o nó com a fonte de maior tensão.
\item eliminar as fontes de corrente.
\item garantir que todas as correntes sejam positivas.
\end{alternativas}
\resposta{(b)}
\resolucao{O terra é uma referência arbitrária de potencial nulo; os potenciais
dos outros nós são medidos em relação a ele. A escolha não altera as
\emph{diferenças} de potencial nem as correntes --- apenas simplifica as
equações.}
\end{questao}

\begin{questao}[1]{}
No circuito da figura, uma fonte de corrente de $\SI{9}{\ampere}$ alimenta dois
resistores em paralelo ligados ao terra. Determine o potencial do nó $V$.

\circuitoq{%
\begin{circuitikz}[scale=1,transform shape,line width=0.8pt,american]
 \draw (0,0) to[isource,l={$\SI{9}{\ampere}$}] (0,2.6)
       -- (1.5,2.6) coordinate(n);
 \node[above,Vcol] at (1.5,2.75) {$V$};
 \draw (n) to[R,l={$\SI{6}{\ohm}$}] (1.5,0) -- (0,0);
 \draw (n) -- (3.5,2.6) to[R,l={$\SI{3}{\ohm}$}] (3.5,0) -- (1.5,0);
 \node[ground] at (0,0){};
\end{circuitikz}}

\resposta{$V = \SI{18}{\volt}$}
\resolucao{A fonte injeta \SI{9}{\ampere} no nó, que escoam pelos dois resistores
em paralelo:
\[
9 = \frac{V}{6}+\frac{V}{3}=\frac{V+2V}{6}=\frac{3V}{6}=\frac{V}{2}
\ \Rightarrow\ V = \SI{18}{\volt}.
\]
Equivale a $V = I\cdot(6\parallel3) = 9\cdot2 = \SI{18}{\volt}$.}
\end{questao}

\blocoaplicacao

\begin{questao}[2]{}
Determine os potenciais nodais $V_1$ e $V_2$ do circuito da figura.

\circuitoq{%
\begin{circuitikz}[scale=1,transform shape,line width=0.8pt,american]
 \draw (0,0) to[isource,l={$\SI{6}{\ampere}$}] (0,2.8) -- (1.5,2.8) coordinate(a);
 \node[above,Vcol] at (1.5,2.95) {$V_1$};
 \draw (a) to[R,l={$\SI{2}{\ohm}$}] (1.5,0) coordinate(g);
 \draw (a) to[R,l={$\SI{3}{\ohm}$}] (4.5,2.8) coordinate(b);
 \node[above,Vcol] at (4.5,2.95) {$V_2$};
 \draw (b) to[R,l={$\SI{6}{\ohm}$}] (4.5,0) -- (g) -- (0,0);
 \node[ground] at (1.5,0){};
\end{circuitikz}}

\resposta{$V_1 \approx \SI{9.82}{\volt}$; $V_2 \approx \SI{6.55}{\volt}$}
\resolucao{Com o terra na base, escrevemos a lei dos nós em cada nó.\\
\textbf{Nó $V_1$} (a fonte injeta \SI{6}{\ampere}):
\[
6=\frac{V_1}{2}+\frac{V_1-V_2}{3}.
\]
\textbf{Nó $V_2$:}
\[
\frac{V_1-V_2}{3}=\frac{V_2}{6}.
\]
Da segunda equação, $2(V_1-V_2)=V_2 \Rightarrow V_1 = \dfrac{3V_2}{2}$.
Substituindo na primeira e resolvendo:
\[
V_1=\frac{108}{11}\approx\SI{9.82}{\volt},
\qquad
V_2=\frac{72}{11}\approx\SI{6.55}{\volt}.
\]
\textbf{Verificação:} corrente da fonte $= \dfrac{V_1}{2}+\dfrac{V_1-V_2}{3}
= \dfrac{108}{22}+\dfrac{36}{33}\approx 4{,}91+1{,}09 = \SI{6}{\ampere}$. Confere.}
\end{questao}

\begin{questao}[2]{}
No nó da figura, uma fonte injeta $\SI{3}{\ampere}$ e outra retira
$\SI{1}{\ampere}$; dois resistores de $\SI{6}{\ohm}$ e $\SI{3}{\ohm}$ ligam o nó ao
terra. Determine o potencial $V$.

\circuitoq{%
\begin{circuitikz}[scale=1,transform shape,line width=0.8pt,american]
 \draw (0,0) to[isource,l={$\SI{3}{\ampere}$}] (0,2.6) -- (1.5,2.6) coordinate(n);
 \node[above,Vcol] at (1.5,2.75){$V$};
 \draw (n) to[R,l={$\SI{6}{\ohm}$}] (1.5,0) -- (0,0);
 \draw (n) -- (3.5,2.6) to[R,l={$\SI{3}{\ohm}$}] (3.5,0) -- (1.5,0);
 \draw (n) -- (5,2.6) to[isource,l={$\SI{1}{\ampere}$},invert] (5,0) -- (3.5,0);
 \node[ground] at (1.5,0){};
\end{circuitikz}}

\resposta{$V = \SI{4}{\volt}$}
\resolucao{A corrente líquida injetada no nó é $3 - 1 = \SI{2}{\ampere}$, que
escoa pelos dois resistores em paralelo:
\[
V = I_{\text{líq}}\cdot(6\parallel3) = 2\cdot2 = \SI{4}{\volt}.
\]
Conferindo pela lei dos nós: $\dfrac{V}{6}+\dfrac{V}{3} = \dfrac{4}{6}+\dfrac{4}{3}
= 0{,}67 + 1{,}33 = \SI{2}{\ampere}$, igual à corrente líquida injetada.}
\end{questao}

\blocoaprofundamento

\begin{questao}[3]{}
No circuito da figura, há uma fonte de tensão e uma de corrente. Usando análise
nodal, determine o potencial $V_A$.

\circuitoq{%
\begin{circuitikz}[scale=1,transform shape,line width=0.8pt,american]
 \draw (0,0) to[battery1,l_={$\SI{12}{\volt}$}] (0,3) -- (1.5,3)
       to[R,l={$\SI{4}{\ohm}$}] (3.5,3) coordinate(a);
 \node[above,Vcol] at (3.5,3.15) {$V_A$};
 \draw (a) to[R,l={$\SI{6}{\ohm}$}] (3.5,0) coordinate(g) -- (0,0);
 \draw (a) -- (5.5,3);
 \draw (5.5,3) to[isource,l={$\SI{2}{\ampere}$}] (5.5,0) -- (g);
 \node[ground] at (3.5,0){};
\end{circuitikz}}

\resposta{$V_A = \SI{12}{\volt}$}
\resolucao{Chamando de $V_A$ o potencial do nó central e notando que o terminal
superior da fonte de tensão está fixo em \SI{12}{\volt}, a lei dos nós em $A$
(correntes saindo $=0$; a fonte de corrente injeta \SI{2}{\ampere} \emph{no} nó):
\[
\frac{V_A-12}{4}+\frac{V_A}{6}-2=0.
\]
Multiplicando por 12: $3(V_A-12)+2V_A-24=0 \Rightarrow 5V_A = 60
\Rightarrow V_A = \SI{12}{\volt}$.\\
\textbf{Interpretação:} nesse ponto, a corrente que a fonte de tensão empurraria
pelo resistor de \SI{4}{\ohm} é nula ($V_A = \SI{12}{\volt}$ iguala os dois lados),
e toda a corrente de \SI{2}{\ampere} da fonte de corrente escoa pelo resistor de
\SI{6}{\ohm}: $2 = 12/6$. Consistente.}
\end{questao}

\begin{questao}[3]{}
Determine os potenciais nodais $V_1$, $V_2$ e $V_3$ do circuito da figura,
montando e resolvendo o sistema de equações.

\circuitoq{%
\begin{circuitikz}[scale=1,transform shape,line width=0.8pt]
 \draw (0,0) to[\fonteV,l={$\SI{12}{\volt}$}] (0,3) -- (1.4,3)
       to[R,l={$\SI{2}{\ohm}$}] (3.2,3) coordinate(n1);
 \node[above,Vcol] at (3.2,3.2){$V_1$}; \fill (n1) circle (0.07);
 \draw (n1) to[R,l={$\SI{4}{\ohm}$}] (6,3) coordinate(n2);
 \node[above,Vcol] at (6,3.2){$V_2$}; \fill (n2) circle (0.07);
 \draw (n1) to[R,l_={$\SI{4}{\ohm}$}] (3.2,0) coordinate(n3);
 \node[below,Vcol] at (3.2,-0.25){$V_3$}; \fill (n3) circle (0.07);
 \draw (n2) to[R,l={$\SI{8}{\ohm}$}] (6,1.5) coordinate(m) -- (6,0) coordinate(b2);
 \draw (n3) to[R,l_={$\SI{4}{\ohm}$}] (0,0);
 \draw (n3) -- (b2);
 \draw (m) -- (7.6,1.5) to[R,l={$\SI{8}{\ohm}$}] (7.6,-1.2) -- (3.2,-1.2) -- (n3);
 \node[ground] at (3.2,-1.2){};
\end{circuitikz}}

\resposta{$V_1 = \dfrac{76}{9}\approx\SI{8.44}{\volt}$,
$V_2=\dfrac{16}{3}\approx\SI{5.33}{\volt}$,
$V_3=\dfrac{40}{9}\approx\SI{4.44}{\volt}$}
\resolucao{\textbf{Montagem.} Com o terra na base e aplicando a lei dos nós
(soma das correntes que \emph{saem} igual a zero) em cada nó:
\[
\text{Nó 1: }\ \frac{V_1-12}{2}+\frac{V_1-V_2}{4}+\frac{V_1-V_3}{4}=0,
\]
\[
\text{Nó 2: }\ \frac{V_2-V_1}{4}+\frac{V_2}{8}+\frac{V_2-V_3}{8}=0,
\]
\[
\text{Nó 3: }\ \frac{V_3-V_1}{4}+\frac{V_3-V_2}{8}+\frac{V_3}{4}=0.
\]
\textbf{Forma matricial} (multiplicando para eliminar denominadores e agrupando
as condutâncias):
\[
\begin{bmatrix}
1{,}0 & -0{,}25 & -0{,}25\\
-0{,}25 & 0{,}5 & -0{,}125\\
-0{,}25 & -0{,}125 & 0{,}625
\end{bmatrix}
\begin{bmatrix}V_1\\V_2\\V_3\end{bmatrix}
=
\begin{bmatrix}6\\0\\0\end{bmatrix}
\]
Note a \textbf{simetria} da matriz de condutâncias --- verificação de que a
montagem está correta.

\textbf{Solução:}
\[
V_1=\frac{76}{9}\approx\SI{8.44}{\volt},\quad
V_2=\frac{16}{3}\approx\SI{5.33}{\volt},\quad
V_3=\frac{40}{9}\approx\SI{4.44}{\volt}.
\]
\textbf{Verificação no nó 2:}
$\dfrac{5{,}33-8{,}44}{4}+\dfrac{5{,}33}{8}+\dfrac{5{,}33-4{,}44}{8}
= -0{,}778+0{,}667+0{,}111 = 0$. \checkmark\\
\textbf{Observação estrutural:} a matriz nodal é o \emph{dual} da matriz de
malhas --- diagonal com a soma das condutâncias do nó, fora da diagonal o negativo
da condutância compartilhada.}
\end{questao}

\begin{questao}[3]{}
No circuito da figura, uma fonte de $\SI{6}{\volt}$ está \emph{flutuante} entre os
nós 1 e 2 (nenhum deles é o terra). Uma fonte de corrente de $\SI{2}{\ampere}$
injeta corrente no nó 1. Determine $V_1$ e $V_2$ pelo método do \textbf{supernó}.

\circuitoq{%
\begin{circuitikz}[scale=1,transform shape,line width=0.8pt,american]
 % fonte de corrente
 \draw (0,0) to[isource,l={$\SI{2}{\ampere}$}] (0,3) -- (1.6,3) coordinate(n1);
 \node[above,Vcol] at (1.6,3.2){$V_1$};
 \fill (n1) circle (0.07);
 % R1 do no1 ao terra
 \draw (n1) to[R,l={$R_1=\SI{6}{\ohm}$}] (1.6,0) coordinate(g);
 % fonte flutuante entre no1 e no2
 \draw (n1) to[battery1,l={$\SI{6}{\volt}$}] (4.6,3) coordinate(n2);
 \node[above,Vcol] at (4.6,3.2){$V_2$};
 \fill (n2) circle (0.07);
 % R2 do no2 ao terra
 \draw (n2) to[R,l={$R_2=\SI{3}{\ohm}$}] (4.6,0) -- (g) -- (0,0);
 % marcacao do superno
 \draw[dashed,laranja,line width=1pt,rounded corners=8pt]
       (1.1,2.4) rectangle (5.1,3.75);
 \node[laranja,font=\footnotesize\sffamily,right] at (5.2,3.05) {supernó};
 \node[ground] at (1.6,0){};
\end{circuitikz}}

\resposta{$V_1 = \SI{8}{\volt}$; $V_2 = \SI{2}{\volt}$}
\resolucao{\textbf{Por que o método normal falha:} a corrente que atravessa a
fonte de \SI{6}{\volt} é desconhecida, então não se pode escrever a lei dos nós
isoladamente em 1 ou em 2.

\medskip
\textbf{Equação 1 --- lei dos nós no supernó.} Traçamos uma superfície fechada
englobando os nós 1 e 2 (tracejado laranja). Tudo o que \emph{entra} deve
\emph{sair}. Entra \SI{2}{\ampere} da fonte de corrente; saem as correntes por
$R_1$ e $R_2$:
\[
\frac{V_1}{R_1}+\frac{V_2}{R_2} = 2
\quad\Longrightarrow\quad
\frac{V_1}{6}+\frac{V_2}{3} = 2.
\]
Note que a corrente \emph{interna} à fonte não aparece --- ela entra e sai da
mesma superfície, cancelando-se. \textbf{Este é o truque do método.}

\medskip
\textbf{Equação 2 --- restrição da fonte.} A fonte impõe diretamente:
\[
V_1 - V_2 = 6.
\]

\medskip
\textbf{Resolvendo.} Da segunda, $V_2 = V_1 - 6$. Substituindo na primeira e
multiplicando por 6:
\[
V_1 + 2(V_1-6) = 12
\;\Rightarrow\;
3V_1 = 24
\;\Rightarrow\;
V_1 = \SI{8}{\volt},\quad V_2 = \SI{2}{\volt}.
\]
\textbf{Verificação:} $\dfrac{8}{6}+\dfrac{2}{3} = 1{,}33+0{,}67 =
\SI{2}{\ampere}$. Confere com a fonte de corrente.}
\end{questao}

% BEGIN BANCO COMPLEMENTAR Análise Nodal
% =====================================================================
% Banco complementar --- Analise Nodal (versao visual revisada)
% N1: fundamentos; N2/N3/N4: bancos separados, todos com circuitos.
% =====================================================================

\subsubsection*{Banco complementar --- Análise Nodal}

\begin{atencao}[Organização do banco revisado]
Nos níveis 2, 3 e 4, o enunciado parte de um \textbf{circuito desenhado}, não de
uma descrição textual da topologia. As questões avançam de redes de um e dois
nós para supernós, fontes dependentes, análise nodal modificada, projeto,
sensibilidade e determinação de equivalentes. O objetivo é treinar a leitura do
circuito e a montagem das equações, e não apenas manipular sistemas já prontos.
\end{atencao}

\blocofixacao

\begin{questao}[1]{}
Em um circuito com cinco nós, quantas tensões nodais independentes existem após a escolha do nó de referência?
\resposta{4}
\resolucao{Com um nó fixado em zero, restam $5-1=4$ potenciais desconhecidos.}
\end{questao}

\begin{questao}[1]{}
Um nó $A$ recebe $\SI{3}{\ampere}$ e está ligado ao terra por $\SI{10}{\ohm}$ e $\SI{15}{\ohm}$. Determine $V_A$.
\resposta{$V_A=\SI{18}{\volt}$}
\resolucao{$V_A(1/10+1/15)=3\Rightarrow V_A=\SI{18}{\volt}$.}
\end{questao}

\begin{questao}[1]{}
Escreva a corrente, no sentido $A\to B$, em um resistor $R$ entre os nós $A$ e $B$.
\resposta{$i_{AB}=(V_A-V_B)/R$}
\resolucao{É a lei de Ohm escrita diretamente em função das tensões nodais.}
\end{questao}

\begin{questao}[1]{}
Uma fonte de corrente ideal de $I$ sai de um nó. Na forma $\mathbf G\mathbf V=\mathbf I$, qual é o sinal de sua contribuição no vetor de correntes injetadas?
\resposta{Negativo}
\resolucao{O vetor $\mathbf I$ representa corrente líquida injetada. Uma fonte que retira corrente entra como $-I$.}
\end{questao}

\begin{questao}[1]{}
Converta $\SI{4}{\ohm}$, $\SI{5}{\ohm}$ e $\SI{20}{\ohm}$, todos ligados do mesmo nó ao terra, em uma única condutância equivalente.
\resposta{$G_{eq}=\SI{0.5}{\siemens}$}
\resolucao{$G_{eq}=1/4+1/5+1/20=0{,}5\,\siemens$.}
\end{questao}

\begin{questao}[1]{}
Uma fonte ideal de tensão de $\SI{12}{\volt}$ tem o terminal negativo aterrado. Qual é a tensão do terminal positivo?
\resposta{$\SI{12}{\volt}$}
\resolucao{O terminal positivo é um nó de potencial conhecido; não é necessário escrever uma equação nodal para determiná-lo.}
\end{questao}

\begin{questao}[1]{}
Quando uma fonte ideal de tensão liga dois nós desconhecidos e nenhum deles é o terra, qual técnica nodal é indicada?
\resposta{Supernó}
\resolucao{A corrente interna da fonte não é conhecida por uma relação resistiva; escreve-se a LKC no contorno do supernó e acrescenta-se a restrição de tensão.}
\end{questao}

\begin{questao}[1]{}
Após obter $V_A=\SI{18}{\volt}$, determine a corrente em um resistor de $\SI{15}{\ohm}$ entre $A$ e o terra.
\resposta{$\SI{1.2}{\ampere}$, de $A$ para o terra}
\resolucao{$i=18/15=\SI{1.2}{\ampere}$.}
\end{questao}

% =====================================================================
% NIVEL 2 --- ANALISE NODAL: APLICACAO COM CIRCUITOS
% 13 questoes, todas com figura
% =====================================================================
\blocoaplicacao

\begin{questao}[2]{}
Determine $V$ e as correntes nos dois resistores.
\circuitoq{\begin{circuitikz}[american,scale=.95,transform shape]
\draw (0,0) to[isource,l_={$\SI{5}{\ampere}$}] (0,2.7) -- (1.8,2.7) coordinate(n);
\draw (n) to[R,l={$\SI{10}{\ohm}$}] (1.8,0) -- (0,0);
\draw (n) -- (3.8,2.7) to[R,l={$\SI{20}{\ohm}$}] (3.8,0) -- (1.8,0);
\node[above,Vcol] at (n){$V$}; \node[ground] at (1.8,0){};
\end{circuitikz}}
\resposta{$V=\SI{33.33}{\volt}$; $i_{10}=\SI{3.33}{\ampere}$; $i_{20}=\SI{1.67}{\ampere}$}
\resolucao{$5=V/10+V/20\Rightarrow V=100/3\,\si{\volt}$. As correntes são $V/R$.}
\end{questao}

\begin{questao}[2]{}
Determine a tensão nodal $V$.
\circuitoq{\begin{circuitikz}[american,scale=.9,transform shape]
\draw (0,0) to[isource,l_={$\SI{6}{\ampere}$}] (0,2.6) -- (1.4,2.6) coordinate(n);
\draw (n) to[R,l={$\SI{8}{\ohm}$}] (1.4,0) -- (0,0);
\draw (n) -- (3.3,2.6) to[R,l={$\SI{8}{\ohm}$}] (3.3,0) -- (1.4,0);
\draw (n) -- (5,2.6) to[isource,l={$\SI{2}{\ampere}$},invert] (5,0) -- (3.3,0);
\node[above,Vcol] at (n){$V$}; \node[ground] at (1.4,0){};
\end{circuitikz}}
\resposta{$V=\SI{16}{\volt}$}
\resolucao{A corrente líquida injetada é $6-2=4$ A. Logo $V(1/8+1/8)=4\Rightarrow V=16$ V.}
\end{questao}

\begin{questao}[2]{}
Determine $V_A$ usando diretamente a LKC no nó.
\circuitoq{\begin{circuitikz}[american,scale=.9,transform shape]
\draw (0,0) to[battery1,l_={$\SI{12}{\volt}$}] (0,2.7) to[R,l={$\SI{4}{\ohm}$}] (2.1,2.7) coordinate(a);
\draw (a) to[R,l={$\SI{12}{\ohm}$}] (2.1,0) -- (0,0);
\draw (a) to[R,l={$\SI{6}{\ohm}$}] (4.4,2.7) -- (4.4,0) to[battery1,l={$\SI{6}{\volt}$},invert] (4.4,2.7);
\draw (4.4,0) -- (2.1,0); \node[ground] at (2.1,0){}; \node[above,Vcol] at (a){$V_A$};
\end{circuitikz}}
\resposta{$V_A=\SI{8}{\volt}$}
\resolucao{$\frac{V_A-12}{4}+\frac{V_A-6}{6}+\frac{V_A}{12}=0\Rightarrow V_A=8$ V.}
\end{questao}

\begin{questao}[2]{}
Determine $V_1$ e $V_2$.
\circuitoq{\begin{circuitikz}[american,scale=.95,transform shape]
\draw (0,0) to[isource,l_={$\SI{3}{\ampere}$}] (0,2.8) -- (1.5,2.8) coordinate(n1);
\draw (n1) to[R,l={$\SI{10}{\ohm}$}] (1.5,0);
\draw (n1) to[R,l={$\SI{10}{\ohm}$}] (4.2,2.8) coordinate(n2);
\draw (n2) to[R,l={$\SI{10}{\ohm}$}] (4.2,0) -- (1.5,0) -- (0,0);
\node[above,Vcol] at (n1){$V_1$}; \node[above,Vcol] at (n2){$V_2$}; \node[ground] at (1.5,0){};
\end{circuitikz}}
\resposta{$V_1=\SI{20}{\volt}$; $V_2=\SI{10}{\volt}$}
\resolucao{$V_1/10+(V_1-V_2)/10=3$ e $V_2/10+(V_2-V_1)/10=0$. Da segunda, $V_1=2V_2$; então $V_2=10$ V.}
\end{questao}

\begin{questao}[2]{}
Determine $V_A$. A fonte e o resistor de $\SI{10}{\ohm}$ formam uma fonte real de tensão.
\circuitoq{\begin{circuitikz}[american,scale=.95,transform shape]
\draw (0,0) to[battery1,l_={$\SI{30}{\volt}$}] (0,2.7) to[R,l={$\SI{10}{\ohm}$}] (2.7,2.7) coordinate(a);
\draw (a) to[R,l={$\SI{15}{\ohm}$}] (2.7,0) -- (0,0); \node[ground] at (2.7,0){}; \node[above,Vcol] at (a){$V_A$};
\end{circuitikz}}
\resposta{$V_A=\SI{18}{\volt}$}
\resolucao{$\frac{V_A-30}{10}+\frac{V_A}{15}=0\Rightarrow V_A=18$ V.}
\end{questao}

\begin{questao}[2]{}
Determine $V_A$ e a corrente no resistor de $\SI{5}{\ohm}$, positiva no sentido $A\to$ nó de $\SI{20}{\volt}$.
\circuitoq{\begin{circuitikz}[american,scale=.9,transform shape]
\draw (0,0) to[battery1,l_={$\SI{20}{\volt}$}] (0,2.7) to[R,l={$\SI{5}{\ohm}$}] (2.4,2.7) coordinate(a);
\draw (a) to[R,l={$\SI{20}{\ohm}$}] (2.4,0) -- (0,0);
\draw (4.4,0) to[isource,l_={$\SI{2}{\ampere}$}] (4.4,2.7) -- (a); \draw (4.4,0)--(2.4,0);
\node[ground] at (2.4,0){}; \node[above,Vcol] at (a){$V_A$};
\end{circuitikz}}
\resposta{$V_A=\SI{24}{\volt}$; $i=\SI{0.8}{\ampere}$}
\resolucao{$\frac{V_A-20}{5}+\frac{V_A}{20}=2\Rightarrow V_A=24$ V; $i=(24-20)/5=0{,}8$ A.}
\end{questao}

\begin{questao}[2]{}
Determine $V_1$ e $V_2$ na rede em escada.
\circuitoq{\begin{circuitikz}[american,scale=.95,transform shape]
\draw (0,0) to[isource,l_={$\SI{4}{\ampere}$}] (0,2.8) -- (1.4,2.8) coordinate(n1);
\draw (n1) to[R,l={$\SI{5}{\ohm}$}] (1.4,0);
\draw (n1) to[R,l={$\SI{10}{\ohm}$}] (4.1,2.8) coordinate(n2);
\draw (n2) to[R,l={$\SI{10}{\ohm}$}] (4.1,0) -- (1.4,0) -- (0,0);
\node[ground] at (1.4,0){}; \node[above,Vcol] at (n1){$V_1$}; \node[above,Vcol] at (n2){$V_2$};
\end{circuitikz}}
\resposta{$V_1=\SI{16}{\volt}$; $V_2=\SI{8}{\volt}$}
\resolucao{$V_1/5+(V_1-V_2)/10=4$ e $V_2/10+(V_2-V_1)/10=0$.}
\end{questao}

\begin{questao}[2]{}
A fonte de $\SI{1}{\ampere}$ transfere corrente do nó 1 para o nó 2. Determine $V_1$ e $V_2$.
\circuitoq{\begin{circuitikz}[american,scale=.9,transform shape]
\draw (0,0) to[isource,l_={$\SI{5}{\ampere}$}] (0,3) -- (1.5,3) coordinate(n1);
\draw (n1) to[R,l={$\SI{6}{\ohm}$}] (1.5,0);
\draw (n1) to[R,l={$\SI{6}{\ohm}$}] (4.1,3) coordinate(n2);
\draw (n2) to[R,l={$\SI{3}{\ohm}$}] (4.1,0) -- (1.5,0)--(0,0);
\draw (n1) to[isource,l={$\SI{1}{\ampere}$}] (n2);
\node[ground] at (1.5,0){}; \node[above,Vcol] at (n1){$V_1$}; \node[above,Vcol] at (n2){$V_2$};
\end{circuitikz}}
\resposta{$V_1=\SI{15.6}{\volt}$; $V_2=\SI{7.2}{\volt}$}
\resolucao{Nó 1: $V_1/6+(V_1-V_2)/6+1=5$. Nó 2: $V_2/3+(V_2-V_1)/6=1$. A solução é $V_1=78/5$ V e $V_2=36/5$ V.}
\end{questao}

\begin{questao}[2]{}
Determine os três potenciais da escada resistiva.
\circuitoq{\begin{circuitikz}[american,scale=.85,transform shape]
\draw (0,0) to[isource,l_={$\SI{4}{\ampere}$}] (0,2.8) -- (1.2,2.8) coordinate(n1);
\draw (n1) to[R,l={$\SI{10}{\ohm}$}] (1.2,0);
\draw (n1) to[R,l={$\SI{10}{\ohm}$}] (3.7,2.8) coordinate(n2);
\draw (n2) to[R,l={$\SI{10}{\ohm}$}] (3.7,0);
\draw (n2) to[R,l={$\SI{10}{\ohm}$}] (6.2,2.8) coordinate(n3);
\draw (n3) to[R,l={$\SI{10}{\ohm}$}] (6.2,0) -- (1.2,0) -- (0,0);
\node[ground] at (1.2,0){}; \node[above,Vcol] at(n1){$V_1$};\node[above,Vcol] at(n2){$V_2$};\node[above,Vcol] at(n3){$V_3$};
\end{circuitikz}}
\resposta{$V_1=\SI{25}{\volt}$; $V_2=\SI{10}{\volt}$; $V_3=\SI{5}{\volt}$}
\resolucao{As três LKC formam uma matriz tridiagonal. Da última, $V_2=2V_3$; da segunda, $V_1=2{,}5V_2$; na primeira obtém-se $V_2=10$ V.}
\end{questao}

\begin{questao}[2]{}
Duas fontes de corrente injetam corrente em nós diferentes. Determine $V_1$ e $V_2$.
\circuitoq{\begin{circuitikz}[american,scale=.9,transform shape]
\draw (0,0) to[isource,l_={$\SI{4}{\ampere}$}] (0,2.8) -- (1.5,2.8) coordinate(n1);
\draw (n1) to[R,l={$\SI{4}{\ohm}$}] (1.5,0);
\draw (n1) to[R,l={$\SI{8}{\ohm}$}] (4.5,2.8) coordinate(n2);
\draw (n2) to[R,l={$\SI{4}{\ohm}$}] (4.5,0);
\draw (6.2,0) to[isource,l_={$\SI{2}{\ampere}$}] (6.2,2.8) -- (n2); \draw (6.2,0)--(1.5,0)--(0,0);
\node[ground] at (1.5,0){};\node[above,Vcol] at(n1){$V_1$};\node[above,Vcol] at(n2){$V_2$};
\end{circuitikz}}
\resposta{$V_1=\SI{14}{\volt}$; $V_2=\SI{10}{\volt}$}
\resolucao{$V_1/4+(V_1-V_2)/8=4$ e $V_2/4+(V_2-V_1)/8=2$.}
\end{questao}

\begin{questao}[2]{}
Determine $V_1$ e $V_2$.
\circuitoq{\begin{circuitikz}[american,scale=.9,transform shape]
\draw (0,0) to[battery1,l_={$\SI{18}{\volt}$}] (0,2.8) to[R,l={$\SI{6}{\ohm}$}] (2.4,2.8) coordinate(n1);
\draw (n1) to[R,l={$\SI{12}{\ohm}$}] (2.4,0);
\draw (n1) to[R,l={$\SI{6}{\ohm}$}] (5.1,2.8) coordinate(n2);
\draw (n2) to[R,l={$\SI{6}{\ohm}$}] (5.1,0);
\draw (6.7,0) to[isource,l_={$\SI{1}{\ampere}$}] (6.7,2.8)--(n2);\draw (6.7,0)--(0,0);
\node[ground] at (2.4,0){};\node[above,Vcol] at(n1){$V_1$};\node[above,Vcol] at(n2){$V_2$};
\end{circuitikz}}
\resposta{$V_1=\SI{10.5}{\volt}$; $V_2=\SI{8.25}{\volt}$}
\resolucao{$\frac{V_1-18}{6}+\frac{V_1}{12}+\frac{V_1-V_2}{6}=0$ e $\frac{V_2}{6}+\frac{V_2-V_1}{6}=1$.}
\end{questao}

\begin{questao}[2]{}
A fonte horizontal força $\SI{2}{\ampere}$ do nó 1 para o nó 2. Determine as tensões nodais.
\circuitoq{\begin{circuitikz}[american,scale=.9,transform shape]
\draw (0,0) to[battery1,l_={$\SI{24}{\volt}$}] (0,2.8) to[R,l={$\SI{6}{\ohm}$}] (2.6,2.8) coordinate(n1);
\draw (n1) to[R,l={$\SI{12}{\ohm}$}] (2.6,0);
\draw (n1) to[isource,l={$\SI{2}{\ampere}$}] (5.2,2.8) coordinate(n2);
\draw (n2) to[R,l={$\SI{6}{\ohm}$}] (5.2,0)--(0,0);
\node[ground] at (2.6,0){};\node[above,Vcol] at(n1){$V_1$};\node[above,Vcol] at(n2){$V_2$};
\end{circuitikz}}
\resposta{$V_1=\SI{8}{\volt}$; $V_2=\SI{12}{\volt}$}
\resolucao{Nó 1: $(V_1-24)/6+V_1/12+2=0\Rightarrow V_1=8$ V. Nó 2: $V_2/6=2\Rightarrow V_2=12$ V.}
\end{questao}

\begin{questao}[2]{}
Determine $V_1$, $V_2$ e a corrente no resistor de $\SI{4}{\ohm}$ entre os nós.
\circuitoq{\begin{circuitikz}[american,scale=.9,transform shape]
\draw (0,0) to[battery1,l_={$\SI{10}{\volt}$}] (0,2.8) to[R,l={$\SI{2}{\ohm}$}] (2.4,2.8) coordinate(n1);
\draw (n1) to[R,l={$\SI{4}{\ohm}$}] (5,2.8) coordinate(n2);
\draw (n1) to[R,l={$\SI{4}{\ohm}$}] (2.4,0);
\draw (n2) to[R,l={$\SI{8}{\ohm}$}] (5,0)--(0,0);
\node[ground] at (2.4,0){};\node[above,Vcol] at(n1){$V_1$};\node[above,Vcol] at(n2){$V_2$};
\end{circuitikz}}
\resposta{$V_1=\SI{6}{\volt}$; $V_2=\SI{4}{\volt}$; $i_{12}=\SI{0.5}{\ampere}$}
\resolucao{$\frac{V_1-10}{2}+\frac{V_1}{4}+\frac{V_1-V_2}{4}=0$ e $\frac{V_2}{8}+\frac{V_2-V_1}{4}=0$. Logo $V_1=6$ V, $V_2=4$ V e $i_{12}=(6-4)/4=0{,}5$ A.}
\end{questao}

% =====================================================================
% NIVEL 3 --- ANALISE NODAL: APROFUNDAMENTO VISUAL
% 12 questoes; inclui 9 problemas com fontes dependentes
% =====================================================================
\blocoaprofundamento

\begin{questao}[3]{}
No circuito, $i_1$ e $i_2$ têm os sentidos indicados. Determine $i_1+i_2$.
\circuitoq{\begin{circuitikz}[american,scale=.92,transform shape]
\draw (0,0) to[isource,l_={$\SI{5}{\ampere}$}] (0,3) -- (1.8,3) coordinate(n);
\draw (n) to[R,l={$\SI{5}{\ohm}$},i>^={$i_2$}] (1.8,0) -- (0,0);
\draw (n) -- (3.9,3) to[cisource,l={$4i_2$}] (3.9,0) -- (1.8,0);
\draw (n) -- (6,3) to[R,l={$\SI{1}{\ohm}$},i>^={$i_1$}] (6,0) -- (3.9,0);
\node[ground] at (1.8,0){};
\end{circuitikz}}
\resposta{$i_1+i_2=\SI{15}{\ampere}$}
\resolucao{Como os resistores estão em paralelo, $i_1=V/1$ e $i_2=V/5$, logo $i_1=5i_2$. Pela LKC no nó superior, $5+4i_2=i_1+i_2$. Assim $5+4i_2=6i_2$, de onde $i_2=2{,}5$ A, $i_1=12{,}5$ A e $i_1+i_2=15$ A.}
\end{questao}

\begin{questao}[3]{}
A fonte dependente drena $0{,}1V_1$ ampère do nó 2. Determine $V_1$ e $V_2$.
\circuitoq{\begin{circuitikz}[american,scale=.9,transform shape]
\draw (0,0) to[isource,l_={$\SI{6}{\ampere}$}] (0,3) -- (1.5,3) coordinate(n1);
\draw (n1) to[R,l={$\SI{5}{\ohm}$}] (1.5,0);
\draw (n1) to[R,l={$\SI{5}{\ohm}$}] (4.4,3) coordinate(n2);
\draw (n2) to[R,l={$\SI{10}{\ohm}$}] (4.4,0);
\draw (n2) -- (6.3,3) to[cisource,l={$0{,}1V_1$},invert] (6.3,0) -- (0,0);
\node[ground] at (1.5,0){};\node[above,Vcol] at(n1){$V_1$};\node[above,Vcol] at(n2){$V_2$};
\end{circuitikz}}
\resposta{$V_1=\SI{18}{\volt}$; $V_2=\SI{6}{\volt}$}
\resolucao{$V_1/5+(V_1-V_2)/5=6$ e $V_2/10+(V_2-V_1)/5+0{,}1V_1=0$. A solução é $(18,6)$ V.}
\end{questao}

\begin{questao}[3]{}
A fonte de tensão controlada impõe $V_2=2V_1$. Determine $V_1$ e $V_2$.
\circuitoq{\begin{circuitikz}[american,scale=.9,transform shape]
\draw (0,0) to[isource,l_={$\SI{3}{\ampere}$}] (0,3) -- (1.6,3) coordinate(n1);
\draw (n1) to[R,l={$\SI{4}{\ohm}$}] (1.6,0);
\draw (n1) to[R,l={$\SI{8}{\ohm}$}] (4.6,3) coordinate(n2);
\draw (n2) to[cvsource,l={$2V_1$}] (4.6,0) -- (0,0);
\node[ground] at (1.6,0){};\node[above,Vcol] at(n1){$V_1$};\node[above,Vcol] at(n2){$V_2$};
\end{circuitikz}}
\resposta{$V_1=\SI{24}{\volt}$; $V_2=\SI{48}{\volt}$}
\resolucao{A LKC em 1 dá $V_1/4+(V_1-V_2)/8=3$. Com $V_2=2V_1$, resulta $V_1/8=3$, portanto $V_1=24$ V e $V_2=48$ V.}
\end{questao}

\begin{questao}[3]{}
A corrente de controle é $i_x=V_1/4$. A fonte dependente injeta $0{,}5i_x$ no nó 2. Determine $V_1$ e $V_2$.
\circuitoq{\begin{circuitikz}[american,scale=.9,transform shape]
\draw (0,0) to[isource,l_={$\SI{4}{\ampere}$}] (0,3) -- (1.5,3) coordinate(n1);
\draw (n1) to[R,l={$\SI{4}{\ohm}$},i>^={$i_x$}] (1.5,0);
\draw (n1) to[R,l={$\SI{8}{\ohm}$}] (4.5,3) coordinate(n2);
\draw (n2) to[R,l={$\SI{8}{\ohm}$}] (4.5,0);
\draw (6.3,0) to[cisource,l_={$0{,}5i_x$}] (6.3,3)--(n2);\draw (6.3,0)--(0,0);
\node[ground] at (1.5,0){};\node[above,Vcol] at(n1){$V_1$};\node[above,Vcol] at(n2){$V_2$};
\end{circuitikz}}
\resposta{$V_1=\SI{16}{\volt}$; $V_2=\SI{16}{\volt}$}
\resolucao{$V_1/4+(V_1-V_2)/8=4$ e $V_2/8+(V_2-V_1)/8=0{,}5(V_1/4)$. A solução é $V_1=V_2=16$ V; logo o resistor entre os nós não conduz corrente.}
\end{questao}

\begin{questao}[3]{}
A fonte de tensão controlada entre os nós satisfaz $V_1-V_2=5i_x$, com $i_x=V_2/5$. Determine $V_1$ e $V_2$ por supernó.
\circuitoq{\begin{circuitikz}[american,scale=.9,transform shape]
\draw (0,0) to[isource,l_={$\SI{5}{\ampere}$}] (0,3)--(1.5,3) coordinate(n1);
\draw (n1) to[R,l={$\SI{10}{\ohm}$}] (1.5,0);
\draw (n1) to[cvsource,l={$5i_x$}] (4.7,3) coordinate(n2);
\draw (n2) to[R,l={$\SI{5}{\ohm}$},i>^={$i_x$}] (4.7,0)--(0,0);
\node[ground] at (1.5,0){};\node[above,Vcol] at(n1){$V_1$};\node[above,Vcol] at(n2){$V_2$};
\end{circuitikz}}
\resposta{$V_1=\SI{25}{\volt}$; $V_2=\SI{12.5}{\volt}$}
\resolucao{Supernó: $V_1/10+V_2/5=5$. Restrição: $V_1-V_2=5(V_2/5)=V_2$, logo $V_1=2V_2$. Da LKC, $0{,}4V_2=5$, portanto $V_2=12{,}5$ V e $V_1=25$ V.}
\end{questao}

\begin{questao}[3]{}
Determine $V_1$, $V_2$ e a corrente na fonte de $\SI{20}{\volt}$, adotada positiva de 1 para 2.
\circuitoq{\begin{circuitikz}[american,scale=.9,transform shape]
\draw (0,0) to[isource,l_={$\SI{4}{\ampere}$}] (0,3)--(1.5,3) coordinate(n1);
\draw (n1) to[R,l={$\SI{10}{\ohm}$}] (1.5,0);
\draw (n1) to[battery1,l={$\SI{20}{\volt}$}] (4.7,3) coordinate(n2);
\draw (n2) to[R,l={$\SI{10}{\ohm}$}] (4.7,0)--(0,0);
\node[ground] at (1.5,0){};\node[above,Vcol] at(n1){$V_1$};\node[above,Vcol] at(n2){$V_2$};
\end{circuitikz}}
\resposta{$V_1=\SI{30}{\volt}$; $V_2=\SI{10}{\volt}$; $i_s=\SI{1}{\ampere}$ de 1 para 2}
\resolucao{Supernó: $V_1/10+V_2/10=4$ e $V_1-V_2=20$. Logo $(V_1,V_2)=(30,10)$ V. No nó 1, dos $4$ A injetados, $3$ A descem pelo resistor; o $1$ A restante atravessa a fonte de tensão de 1 para 2.}
\end{questao}

\begin{questao}[3]{}
Determine $V_1$, $V_2$, $V_3$ e a corrente $i_{23}$ no resistor de ponte de $\SI{6}{\ohm}$, positiva de 2 para 3.
\circuitoq{\begin{circuitikz}[american,scale=.82,transform shape]
\coordinate (g) at (0,-1.6);
\draw (g) to[isource,l_={$\SI{3}{\ampere}$}] (0,1.6) -- (1.3,1.6) coordinate(n1);
\draw (n1) to[R,l={$\SI{6}{\ohm}$}] (4.2,3.0) coordinate(n2);
\draw (n1) to[R,l_={$\SI{12}{\ohm}$}] (4.2,0.2) coordinate(n3);
\draw (n2) to[R,l={$\SI{6}{\ohm}$}] (n3);
\draw (n2) to[R,l={$\SI{12}{\ohm}$}] (4.2,-1.6);
\draw (n3) to[R,l_={$\SI{12}{\ohm}$}] (4.2,-1.6);
\draw (4.2,-1.6) -- (g);
\node[ground] at(g){};
\node[above,Vcol] at(n1){$V_1$};\node[above,Vcol] at(n2){$V_2$};\node[right,Vcol] at(n3){$V_3$};
\end{circuitikz}}
\resposta{$V_1=\dfrac{576}{19}\,\si{\volt}\approx\SI{30.32}{\volt}$; $V_2=\dfrac{360}{19}\,\si{\volt}\approx\SI{18.95}{\volt}$; $V_3=\dfrac{324}{19}\,\si{\volt}\approx\SI{17.05}{\volt}$; $i_{23}=\dfrac{6}{19}\,\si{\ampere}\approx\SI{0.316}{\ampere}$}
\resolucao{As LKC são $(V_1-V_2)/6+(V_1-V_3)/12=3$, $V_2/12+(V_2-V_1)/6+(V_2-V_3)/6=0$ e $V_3/12+(V_3-V_1)/12+(V_3-V_2)/6=0$. A solução é $(576,360,324)/19$ V. Assim $i_{23}=(V_2-V_3)/6=6/19$ A.}
\end{questao}

\begin{questao}[3]{}
A fonte controlada transfere corrente do nó 1 para o nó 2 no valor $0{,}05V_2$. Determine $V_1$ e $V_2$.
\circuitoq{\begin{circuitikz}[american,scale=.9,transform shape]
\draw (0,0) to[isource,l_={$\SI{5}{\ampere}$}] (0,3)--(1.5,3) coordinate(n1);
\draw (n1) to[R,l={$\SI{5}{\ohm}$}] (1.5,0);
\draw (n1) to[R,l={$\SI{10}{\ohm}$}] (4.5,3) coordinate(n2);
\draw (n2) to[R,l={$\SI{10}{\ohm}$}] (4.5,0);
\draw (n1) to[cisource,l={$0{,}05V_2$}] (n2);\draw (4.5,0)--(0,0);
\node[ground] at(1.5,0){};\node[above,Vcol] at(n1){$V_1$};\node[above,Vcol] at(n2){$V_2$};
\end{circuitikz}}
\resposta{$V_1=\SI{18.75}{\volt}$; $V_2=\SI{12.5}{\volt}$}
\resolucao{$V_1/5+(V_1-V_2)/10+0{,}05V_2=5$ e $V_2/10+(V_2-V_1)/10-0{,}05V_2=0$. Resulta $V_1=75/4$ V e $V_2=25/2$ V.}
\end{questao}

\begin{questao}[3]{}
A fonte flutuante é controlada e impõe $V_1-V_2=0{,}5V_2$. Determine as tensões nodais.
\circuitoq{\begin{circuitikz}[american,scale=.9,transform shape]
\draw (0,0) to[isource,l_={$\SI{3}{\ampere}$}] (0,3)--(1.5,3) coordinate(n1);
\draw (n1) to[R,l={$\SI{6}{\ohm}$}] (1.5,0);
\draw (n1) to[cvsource,l={$0{,}5V_2$}] (4.7,3) coordinate(n2);
\draw (n2) to[R,l={$\SI{3}{\ohm}$}] (4.7,0)--(0,0);
\node[ground] at(1.5,0){};\node[above,Vcol] at(n1){$V_1$};\node[above,Vcol] at(n2){$V_2$};
\end{circuitikz}}
\resposta{$V_1=\dfrac{54}{7}\,\si{\volt}\approx\SI{7.71}{\volt}$; $V_2=\dfrac{36}{7}\,\si{\volt}\approx\SI{5.14}{\volt}$}
\resolucao{Supernó: $V_1/6+V_2/3=3$. A restrição dá $V_1=1{,}5V_2$. Substituindo: $7V_2/12=3\Rightarrow V_2=36/7$ V.}
\end{questao}

\begin{questao}[3]{}
Determine a resistência vista pelo terminal usando a fonte de teste de $\SI{1}{\ampere}$. A fonte dependente injeta $0{,}05V$ no nó.
\circuitoq{\begin{circuitikz}[american,scale=.95,transform shape]
\draw (0,0) to[isource,l_={$I_t=\SI{1}{\ampere}$}] (0,3)--(2,3) coordinate(n);
\draw (n) to[R,l={$\SI{10}{\ohm}$}] (2,0)--(0,0);
\draw (4,0) to[cisource,l_={$0{,}05V$}] (4,3)--(n);\draw (4,0)--(2,0);
\node[ground] at(2,0){};\node[above,Vcol] at(n){$V$};
\end{circuitikz}}
\resposta{$R_{in}=\SI{20}{\ohm}$}
\resolucao{$V/10=1+0{,}05V\Rightarrow 0{,}05V=1\Rightarrow V=20$ V. Como $I_t=1$ A, $R_{in}=V/I_t=20\,\ohm$.}
\end{questao}

\begin{questao}[3]{}
No circuito, a fonte dependente injeta $0{,}1V$ ampère. Determine $V$ e verifique o balanço de potência.
\circuitoq{\begin{circuitikz}[american,scale=.9,transform shape]
\draw (0,0) to[isource,l_={$\SI{2}{\ampere}$}] (0,3)--(1.5,3) coordinate(n);
\draw (n) to[R,l={$\SI{4}{\ohm}$}] (1.5,0);
\draw (n)--(3.6,3) to[R,l={$\SI{8}{\ohm}$}] (3.6,0)--(1.5,0);
\draw (5.4,0) to[cisource,l_={$0{,}1V$}] (5.4,3)--(n);\draw (5.4,0)--(1.5,0);
\node[ground] at(1.5,0){};\node[above,Vcol] at(n){$V$};
\end{circuitikz}}
\resposta{$V=\dfrac{80}{11}\,\si{\volt}\approx\SI{7.27}{\volt}$; balanço fecha em $\SI{19.83}{\watt}$}
\resolucao{$V/4+V/8=2+0{,}1V\Rightarrow V=80/11$ V. Os resistores dissipam $V^2(1/4+1/8)=2400/121\approx19{,}83$ W. As fontes entregam $2V+0{,}1V^2=160/11+640/121=2400/121$ W.}
\end{questao}

\begin{questao}[3]{}
A fonte de $\SI{8}{\volt}$ é flutuante. Determine $V_1$, $V_2$ e $V_3$.
\circuitoq{\begin{circuitikz}[american,scale=.85,transform shape]
\draw (0,0) to[isource,l_={$\SI{5}{\ampere}$}] (0,3)--(1.3,3) coordinate(n1);
\draw (n1) to[R,l={$\SI{4}{\ohm}$}] (1.3,0);
\draw (n1) to[battery1,l={$\SI{8}{\volt}$}] (3.9,3) coordinate(n2);
\draw (n2) to[R,l={$\SI{8}{\ohm}$}] (3.9,0);
\draw (n2) to[R,l={$\SI{4}{\ohm}$}] (6.3,3) coordinate(n3);
\draw (n3) to[R,l={$\SI{4}{\ohm}$}] (6.3,0)--(0,0);
\node[ground] at(1.3,0){};\node[above,Vcol] at(n1){$V_1$};\node[above,Vcol] at(n2){$V_2$};\node[above,Vcol] at(n3){$V_3$};
\end{circuitikz}}
\resposta{$V_1=\SI{14}{\volt}$; $V_2=\SI{6}{\volt}$; $V_3=\SI{3}{\volt}$}
\resolucao{Supernó $1$--$2$: $V_1/4+V_2/8+(V_2-V_3)/4=5$; restrição $V_1-V_2=8$; nó 3: $V_3/4+(V_3-V_2)/4=0$. A solução é $(14,6,3)$ V.}
\end{questao}

% =====================================================================
% NIVEL 4 --- ANALISE NODAL: DESAFIOS COM CIRCUITOS
% 10 questoes, todas com figura
% =====================================================================
\blocodesafio

\begin{questao}[4]{}
A fonte dependente injeta $gV$ no nó. Obtenha $V(g)$, determine o valor crítico de $g$ e calcule $V$ para $g=\SI{0.04}{\siemens}$.
\circuitoq{\begin{circuitikz}[american,scale=.95,transform shape]
\draw (0,0) to[isource,l_={$\SI{1}{\ampere}$}] (0,3)--(2,3) coordinate(n);
\draw (n) to[R,l={$\SI{20}{\ohm}$}] (2,0)--(0,0);
\draw (4.2,0) to[cisource,l_={$gV$}] (4.2,3)--(n);\draw (4.2,0)--(2,0);
\node[ground] at(2,0){};\node[above,Vcol] at(n){$V$};
\end{circuitikz}}
\resposta{$V(g)=1/(0{,}05-g)$; $g_c=\SI{0.05}{\siemens}$; para $g=\SI{0.04}{\siemens}$, $V=\SI{100}{\volt}$}
\resolucao{LKC: $V/20=1+gV$. Assim $V=1/(0{,}05-g)$. Em $g=0{,}05$ S o coeficiente nodal se anula: a rede linear torna-se singular.}
\end{questao}

\begin{questao}[4]{}
A fonte controlada injeta $kV_1$ no nó 2. Determine $k$ para que $V_2=\SI{10}{\volt}$ e encontre $V_1$.
\circuitoq{\begin{circuitikz}[american,scale=.9,transform shape]
\draw (0,0) to[isource,l_={$\SI{2}{\ampere}$}] (0,3)--(1.5,3) coordinate(n1);
\draw (n1) to[R,l={$\SI{5}{\ohm}$}] (1.5,0);
\draw (n1) to[R,l={$\SI{5}{\ohm}$}] (4.5,3) coordinate(n2);
\draw (n2) to[R,l={$\SI{10}{\ohm}$}] (4.5,0);
\draw (6.3,0) to[cisource,l_={$kV_1$}] (6.3,3)--(n2);\draw (6.3,0)--(0,0);
\node[ground] at(1.5,0){};\node[above,Vcol] at(n1){$V_1$};\node[above,Vcol] at(n2){$V_2$};
\end{circuitikz}}
\resposta{$k=\SI{0.1}{\siemens}$; $V_1=\SI{10}{\volt}$}
\resolucao{Impondo $V_2=10$ V na LKC do nó 1: $V_1/5+(V_1-10)/5=2\Rightarrow V_1=10$ V. No nó 2: $10/10+(10-10)/5=k(10)$, logo $k=0{,}1$ S.}
\end{questao}

\begin{questao}[4]{}
A fonte dependente injeta $2i_x$ no nó 3, com $i_x=(V_1-V_2)/4$. Determine as três tensões e identifique um ramo que fica sem corrente.
\circuitoq{\begin{circuitikz}[american,scale=.82,transform shape]
\draw (0,0) to[isource,l_={$\SI{4}{\ampere}$}] (0,3)--(1.3,3) coordinate(n1);
\draw (n1) to[R,l={$\SI{4}{\ohm}$}] (1.3,0);
\draw (n1) to[R,l={$\SI{4}{\ohm}$},i>^={$i_x$}] (3.8,3) coordinate(n2);
\draw (n2) to[R,l={$\SI{8}{\ohm}$}] (3.8,0);
\draw (n2) to[R,l={$\SI{4}{\ohm}$}] (6.3,3) coordinate(n3);
\draw (n3) to[R,l={$\SI{4}{\ohm}$}] (6.3,0);
\draw (7.9,0) to[cisource,l_={$2i_x$}] (7.9,3)--(n3);\draw (7.9,0)--(0,0);
\node[ground] at(1.3,0){};\node[above,Vcol] at(n1){$V_1$};\node[above,Vcol] at(n2){$V_2$};\node[above,Vcol] at(n3){$V_3$};
\end{circuitikz}}
\resposta{$V_1=\SI{12}{\volt}$; $V_2=\SI{8}{\volt}$; $V_3=\SI{8}{\volt}$; o resistor entre 2 e 3 conduz corrente nula}
\resolucao{As LKC são $V_1/4+(V_1-V_2)/4=4$, $V_2/8+(V_2-V_1)/4+(V_2-V_3)/4=0$ e $V_3/4+(V_3-V_2)/4=2(V_1-V_2)/4$. A solução é $(12,8,8)$ V.}
\end{questao}

\begin{questao}[4]{}
A fonte de tensão controlada satisfaz $V_1-V_2=5i_x$, $i_x=V_2/10$, e a fonte de corrente controlada injeta $0{,}05V_1$ no nó 2. Determine $V_1$ e $V_2$.
\circuitoq{\begin{circuitikz}[american,scale=.88,transform shape]
\draw (0,0) to[isource,l_={$\SI{3}{\ampere}$}] (0,3)--(1.4,3) coordinate(n1);
\draw (n1) to[R,l={$\SI{10}{\ohm}$}] (1.4,0);
\draw (n1) to[cvsource,l={$5i_x$}] (4.5,3) coordinate(n2);
\draw (n2) to[R,l={$\SI{10}{\ohm}$},i>^={$i_x$}] (4.5,0);
\draw (6.3,0) to[cisource,l_={$0{,}05V_1$}] (6.3,3)--(n2);\draw (6.3,0)--(0,0);
\node[ground] at(1.4,0){};\node[above,Vcol] at(n1){$V_1$};\node[above,Vcol] at(n2){$V_2$};
\end{circuitikz}}
\resposta{$V_1=\dfrac{180}{7}\,\si{\volt}\approx\SI{25.71}{\volt}$; $V_2=\dfrac{120}{7}\,\si{\volt}\approx\SI{17.14}{\volt}$}
\resolucao{A restrição é $V_1-V_2=5(V_2/10)=0{,}5V_2$, logo $V_1=1{,}5V_2$. No supernó: $V_1/10+V_2/10=3+0{,}05V_1$. Substituindo, $0{,}175V_2=3$.}
\end{questao}

\begin{questao}[4]{}
Determine o equivalente de Thévenin visto no nó 2. A fonte dependente injeta $0{,}05V_1$ no nó 2.
\circuitoq{\begin{circuitikz}[american,scale=.9,transform shape]
\draw (0,0) to[isource,l_={$\SI{2}{\ampere}$}] (0,3)--(1.5,3) coordinate(n1);
\draw (n1) to[R,l={$\SI{10}{\ohm}$}] (1.5,0);
\draw (n1) to[R,l={$\SI{5}{\ohm}$}] (4.7,3) coordinate(n2);
\draw (6.5,0) to[cisource,l_={$0{,}05V_1$}] (6.5,3)--(n2);\draw (6.5,0)--(0,0);
\draw (n2)--(5.7,3);\draw (5.7,2.75)--(5.7,3.25);\node[right] at(5.8,3){porta};
\node[ground] at(1.5,0){};\node[above,Vcol] at(n1){$V_1$};\node[above,Vcol] at(n2){$V_2$};
\end{circuitikz}}
\resposta{$V_{Th}=\SI{50}{\volt}$; $R_{Th}=\SI{30}{\ohm}$}
\resolucao{Com a porta aberta: $(V_2-V_1)/5=0{,}05V_1\Rightarrow V_2=1{,}25V_1$. No nó 1, $V_1/10+(V_1-V_2)/5=2$, resultando $V_1=40$ V e $V_{Th}=50$ V. Para $R_{Th}$, abra a fonte de $2$ A e injete $I_t$ na porta. Da LKC do nó 1, $V_1=(2/3)V_2$; no nó 2 obtém-se $I_t=V_2/30$, logo $R_{Th}=30\,\ohm$.}
\end{questao}

\begin{questao}[4]{}
A fonte controlada impõe $V_2=kV_1$. Obtenha $V_1(k)$ e determine o valor de $k$ para o qual a solução se torna singular.
\circuitoq{\begin{circuitikz}[american,scale=.9,transform shape]
\draw (0,0) to[isource,l_={$\SI{3}{\ampere}$}] (0,3)--(1.5,3) coordinate(n1);
\draw (n1) to[R,l={$\SI{10}{\ohm}$}] (1.5,0);
\draw (n1) to[R,l={$\SI{5}{\ohm}$}] (4.5,3) coordinate(n2);
\draw (n2) to[cvsource,l={$kV_1$}] (4.5,0)--(0,0);
\node[ground] at(1.5,0){};\node[above,Vcol] at(n1){$V_1$};\node[above,Vcol] at(n2){$V_2$};
\end{circuitikz}}
\resposta{$V_1=3/(0{,}3-0{,}2k)$; singular em $k=1{,}5$}
\resolucao{LKC no nó 1: $V_1/10+(V_1-kV_1)/5=3$. O coeficiente é $0{,}1+0{,}2(1-k)=0{,}3-0{,}2k$. Ele zera em $k=1{,}5$.}
\end{questao}

\begin{questao}[4]{}
Escolha $R_L$ para obter $V=\SI{20}{\volt}$. Depois calcule a sensibilidade $dV/dR_L$ no ponto de projeto. A fonte dependente injeta $0{,}02V$ no nó.
\circuitoq{\begin{circuitikz}[american,scale=.88,transform shape]
\draw (0,0) to[isource,l_={$\SI{2}{\ampere}$}] (0,3)--(1.5,3) coordinate(n);
\draw (n) to[R,l={$\SI{20}{\ohm}$}] (1.5,0);
\draw (n)--(3.7,3) to[R,l={$R_L$}] (3.7,0)--(1.5,0);
\draw (5.5,0) to[cisource,l_={$0{,}02V$}] (5.5,3)--(n);\draw (5.5,0)--(1.5,0);
\node[ground] at(1.5,0){};\node[above,Vcol] at(n){$V$};
\end{circuitikz}}
\resposta{$R_L=\dfrac{100}{7}\,\ohm\approx\SI{14.29}{\ohm}$; $dV/dR_L\approx\SI{0.98}{\volt\per\ohm}$}
\resolucao{$V(1/20+1/R_L-0{,}02)=2$, portanto $V=2R_L/(1+0{,}03R_L)$. Impondo $V=20$ V, resulta $R_L=100/7\,\ohm$. A derivada é $dV/dR_L=2/(1+0{,}03R_L)^2\approx0{,}98$ V/$\ohm$.}
\end{questao}

\begin{questao}[4]{}
Resolva o circuito por análise nodal modificada, tomando como incógnitas $V_1$, $V_2$ e a corrente $i_s$ da fonte flutuante de $\SI{10}{\volt}$. A fonte controlada drena $0{,}05V_1$ do nó 2.
\circuitoq{\begin{circuitikz}[american,scale=.88,transform shape]
\draw (0,0) to[isource,l_={$\SI{3}{\ampere}$}] (0,3)--(1.5,3) coordinate(n1);
\draw (n1) to[R,l={$\SI{10}{\ohm}$}] (1.5,0);
\draw (n1) to[battery1,l={$\SI{10}{\volt}$}] (4.7,3) coordinate(n2);
\draw (n2) to[R,l={$\SI{20}{\ohm}$}] (4.7,0);
\draw (n2)--(6.5,3) to[cisource,l={$0{,}05V_1$}] (6.5,0)--(0,0);
\node[ground] at(1.5,0){};\node[above,Vcol] at(n1){$V_1$};\node[above,Vcol] at(n2){$V_2$};
\end{circuitikz}}
\resposta{$V_1=\SI{17.5}{\volt}$; $V_2=\SI{7.5}{\volt}$; $i_s=\SI{1.25}{\ampere}$ de 1 para 2}
\resolucao{Com $i_s$ positivo de 1 para 2: $V_1/10+i_s=3$; $V_2/20+0{,}05V_1-i_s=0$; $V_1-V_2=10$. A solução é $(17{,}5,7{,}5,1{,}25)$.}
\end{questao}

\begin{questao}[4]{}
Determine $R$ para que $V_2=\SI{10}{\volt}$. A fonte dependente drena $0{,}05V_1$ do nó 2.
\circuitoq{\begin{circuitikz}[american,scale=.9,transform shape]
\draw (0,0) to[isource,l_={$\SI{5}{\ampere}$}] (0,3)--(1.5,3) coordinate(n1);
\draw (n1) to[R,l={$\SI{10}{\ohm}$}] (1.5,0);
\draw (n1) to[R,l={$\SI{5}{\ohm}$}] (4.5,3) coordinate(n2);
\draw (n2) to[R,l={$R$}] (4.5,0);
\draw (n2)--(6.3,3) to[cisource,l={$0{,}05V_1$}] (6.3,0)--(0,0);
\node[ground] at(1.5,0){};\node[above,Vcol] at(n1){$V_1$};\node[above,Vcol] at(n2){$V_2$};
\end{circuitikz}}
\resposta{$V_1=\dfrac{70}{3}\,\si{\volt}$; $R=\dfrac{20}{3}\,\ohm\approx\SI{6.67}{\ohm}$}
\resolucao{Com $V_2=10$ V, $V_1/10+(V_1-10)/5=5\Rightarrow V_1=70/3$ V. No nó 2: $10/R+(10-V_1)/5+0{,}05V_1=0$, de onde $10/R=1{,}5$ e $R=20/3\,\ohm$.}
\end{questao}

\begin{questao}[4]{}
Determine $V_1$, $V_2$, $V_3$ e a corrente $i_{23}$ no resistor de ponte, positiva de 2 para 3. A fonte dependente injeta $0{,}125V_2$ no nó 3.
\circuitoq{\begin{circuitikz}[american,scale=.82,transform shape]
\coordinate (g) at (0,-1.7);
\draw (g) to[isource,l_={$\SI{6}{\ampere}$}] (0,1.6) -- (1.2,1.6) coordinate(n1);
\draw (n1) to[R,l={$\SI{4}{\ohm}$}] (4.0,3.0) coordinate(n2);
\draw (n1) to[R,l_={$\SI{4}{\ohm}$}] (4.0,0.2) coordinate(n3);
\draw (n2) to[R,l={$\SI{4}{\ohm}$}] (n3);
\draw (n2) to[R,l={$\SI{8}{\ohm}$}] (6.2,3.0) -- (6.2,-1.7) -- (g);
\draw (n3) to[R,l_={$\SI{8}{\ohm}$}] (4.0,-1.7) -- (g);
\draw (7.8,-1.7) to[cisource,l_={$0{,}125V_2$}] (7.8,0.2) -- (n3);
\draw (7.8,-1.7) -- (g);
\node[ground] at(g){};
\node[above,Vcol] at(n1){$V_1$};\node[above,Vcol] at(n2){$V_2$};\node[right,Vcol] at(n3){$V_3$};
\end{circuitikz}}
\resposta{$V_1=\SI{57}{\volt}$; $V_2=\SI{42}{\volt}$; $V_3=\SI{48}{\volt}$; $i_{23}=\SI{-1.5}{\ampere}$}
\resolucao{As LKC são $(V_1-V_2)/4+(V_1-V_3)/4=6$, $V_2/8+(V_2-V_1)/4+(V_2-V_3)/4=0$ e $V_3/8+(V_3-V_1)/4+(V_3-V_2)/4=0{,}125V_2$. A solução é $(57,42,48)$ V. Assim $i_{23}=(42-48)/4=-1{,}5$ A, isto é, $1{,}5$ A de 3 para 2.}
\end{questao}


% END BANCO COMPLEMENTAR

\gabarito[Capítulo 7 --- Análise Nodal]
